\documentclass[12pt,a4paper]{ctexart}
\usepackage[margin=2.4cm,top=2.0cm,bottom=2.0cm]{geometry}
\usepackage{array}
\setCJKmainfont{SimSun}[AutoFakeBold=2.5]
\pagestyle{empty}
\setlength{\parindent}{0pt}
\begin{document}
\noindent\rule{\textwidth}{0.6pt}

\vspace{0.12cm}
\hfill
\begin{tabular}{|c|c|}
\hline
\rule{0pt}{15pt}\hspace{4pt}得\ \ 分\hspace{4pt} & \rule{0pt}{15pt}\hspace{3.0cm} \\
\hline
\end{tabular}

\vspace{2.5cm}
\begin{center}
{\zihao{2}\bfseries 《数值分析实践》综合实验报告}\\[0.5cm]
{\zihao{4} （2025-2026-2）}
\end{center}

\vspace{1.3cm}
\begin{center}
\renewcommand{\arraystretch}{1.9}
\begin{tabular}{|c|p{9.2cm}|}
\hline
实验名称1 & 个性化非线性方程迭代求解及AI生成代码的批判性验证 \\
\hline
实验名称2 & 线性方程组直接解法与迭代解法的实现、对比及应用 \\
\hline
\end{tabular}
\end{center}

\vspace{3.1cm}
\begin{center}
\renewcommand{\arraystretch}{2.15}
\begin{tabular}{r@{\ }l}
班\qquad 级： & \underline{\makebox[5.4cm][c]{数学与应用数学}} \\[0.15cm]
学\qquad 号： & \underline{\makebox[5.4cm][c]{}} \\[0.15cm]
姓\qquad 名： & \underline{\makebox[5.4cm][c]{VesperaZephyr}} \\[0.15cm]
指导教师： & \underline{\makebox[5.4cm][c]{}} \\
\end{tabular}
\end{center}

\vfill
\begin{center}
{\zihao{4} 2026 年 6 月 8 日}
\end{center}
\end{document}
